% 相邻样本经同一算法得到相近输出，从而在任意测试点上损失接近。
\begin{tikzpicture}[figure picture,foundation picture,
  data node/.style={foundation node,fill=FoundationInput!10,minimum width=27mm,minimum height=12mm,inner sep=3pt},
  alg node/.style={foundation node,fill=FoundationModel!10,minimum width=21mm,minimum height=10mm},
  out node/.style={foundation node,fill=FoundationOutput!8,minimum width=25mm,minimum height=10mm},
  note/.style={font=\figurefont\footnotesize,text=FigureMuted,align=center}
]
  \node[data node] (s) {$S=(z_1,\ldots,z_i,\ldots,z_n)$};
  \node[data node,below=14mm of s] (sp) {$S'=(z_1,\ldots,z_i',\ldots,z_n)$};
  \draw[FoundationLoss,line width=1.1pt,<->] (s) -- node[right,note,text=FoundationLoss]{只替换\\一个观测} (sp);

  \node[alg node,right=15mm of s] (a1) {同一算法$A$};
  \node[alg node,right=15mm of sp] (a2) {同一算法$A$};
  \node[out node,right=15mm of a1] (w1) {$w=A(S)$};
  \node[out node,right=15mm of a2] (w2) {$w'=A(S')$};
  \draw[foundation flow] (s)--(a1);
  \draw[foundation flow] (sp)--(a2);
  \draw[foundation flow] (a1)--(w1);
  \draw[foundation flow] (a2)--(w2);
  \draw[FoundationInput,line width=1.1pt,<->] (w1)--node[right,note,text=FoundationInput]{损失差}(w2);

  \node[foundation node,fill=FoundationLoss!9,below right=4mm and 12mm of w1] (test)
    {$\forall z:\quad|\ell(w,z)-\ell(w',z)|\leq\beta_n$};
  \draw[foundation flow,dashed] (w1)--(test);
  \draw[foundation flow,dashed] (w2)--(test);
\end{tikzpicture}%
